\documentclass[../main.tex]{subfiles}

\begin{document}

\justifying

\section{Factor de dos y fuerza en la partícula 2}

\subsection{Directory: 2026-03-03-101553}

La siguiente grafica muestra:
\begin{gather*}
    \frac{1}{2}F_{\mathrm{tabla~swap}} - F_{\mathrm{LAMMPS}}
\end{gather*}

\begin{figure}[!ht]
    \centering
    \includegraphics[width=\textwidth]{../Figures/2026-03-03-101553/fig_FswapDiff.png}
    \caption{
        Línea sólida: Componente $x$, línea punteada: Componente $y$.
        Línea azul: particula central.
        Línea naranja: partícula superior.
        Línea verde: partícula inferior.
    }\label{fig:Config1-SwapDiff}
\end{figure}

Podemos ver que en LAMMPS se registra la mitad del valor que se tabula.
Pero también se observa que aparece una fuerza en la partícula dos.
Esa fuerza no debe de aparecer ahí.
Mi intución es que, la aparición de esa fuerza, ocasiona el ruido de la fuerza en la partícula 1 y 3.
No me queda claro cómo aparece esa fuerza en la partícula 2.

En la siguiente figura se muestra el cálculo de fuerzas usando las distancias y la fuerza de la tabla,
\begin{gather}
    \begin{pmatrix}
        \vec{f}_i \\ \vec{f}_j \\\vec{f}_k 
    \end{pmatrix}
    =
    \begin{pmatrix}
        f_{i1} & f_{i2} & 0 \\
        f_{j1} & 0 & f_{j2} \\
        0 & f_{k1} & f_{k2}
    \end{pmatrix}
    \begin{pmatrix}
        \vec{r}_{ij} \\ \vec{r}_{ik} \\\vec{r}_{jk} 
    \end{pmatrix}
    \label{eqn:lammps3body}
\end{gather}

\newpage

\begin{figure}[!ht]
    \centering
    \includegraphics[width=\textwidth]{../Figures/2026-03-03-101553/fig_Fswaptable.png}
    \caption{
        Línea sólida: Componente $x$, línea punteada: Componente $y$.
        Línea azul: particula central.
        Línea naranja: partícula superior.
        Línea verde: partícula inferior.
    }\label{fig:Config1-SwapDiff}
\end{figure}

Revisando el dump, aún aparece que la partícula 2 tiene una velocidad.
Posiblemente de ahí provenga el error.
Modificaré el script de lammps para que después de que se deje de mover la partícula 2, está no tenga velocidad.

\newpage

\subsection{Directory: 2026-03-03-115610}

Pues no era eso\ldots

\begin{figure}[!ht]
    \centering
    \includegraphics[width=\textwidth]{../Figures/2026-03-03-115610/fig_FswapDiff.png}
    \caption{
        Línea sólida: Componente $x$, línea punteada: Componente $y$.
        Línea azul: particula central.
        Línea naranja: partícula superior.
        Línea verde: partícula inferior.
    }\label{fig:Config1-SwapDiff}
\end{figure}

Voy a agregar un factor de \(2\) en la tabulación de la tabla y hacer el mismo análisis a ver que pasa.

\newpage

\subsection{Directory: 2026-03-03-120832}

\begin{figure}[!ht]
    \centering
    \includegraphics[width=\textwidth]{../Figures/2026-03-03-120832/fig_FswapDiff.png}
    \caption{
        Línea sólida: Componente $x$, línea punteada: Componente $y$.
        Línea azul: particula central.
        Línea naranja: partícula superior.
        Línea verde: partícula inferior.
    }\label{fig:Config1-SwapDiff}
\end{figure}

\newpage
Pondré que la energía sea cero en todo el potencial, pero la fuerza que no sea nula.
Pffff, lo mismo.

\subsection{Directory: 2026-03-03-125817}

\begin{figure}[!ht]
    \centering
    \includegraphics[width=\textwidth]{../Figures/2026-03-03-125817/fig_FswapDiff.png}
    \caption{
        Línea sólida: Componente $x$, línea punteada: Componente $y$.
        Línea azul: particula central.
        Línea naranja: partícula superior.
        Línea verde: partícula inferior.
    }\label{fig:Config1-SwapDiff}
\end{figure}


\end{document}
